\chapter{Query di estrazione del catalogo}
\label{app:query}

La query seguente documenta la selezione iniziale effettuata tramite
CasJobs. I nomi delle colonne esportate coincidono con quelli utilizzati
dalla pipeline locale. Il vincolo di etichetta univoca descritto nella
\cref{sec:etichette} viene applicato dopo l'esportazione.

\begin{minted}{sql}
SELECT
    p.objid,
    p.ra,
    p.dec,
    z.p_el_debiased AS prob_elliptical,
    z.p_cs_debiased AS prob_spiral,
    z.p_mg AS prob_merger,
    s.z AS redshift
INTO mydb.GalaxyCandidates_Final
FROM PhotoObj AS p
JOIN zooSpec AS z ON p.objid = z.objid
JOIN SpecObj AS s ON p.objid = s.bestObjID
WHERE
    p.clean = 1
    AND p.type = 3
    AND p.g BETWEEN 14 AND 17.5
    AND s.z BETWEEN 0.005 AND 0.1
    AND (
        z.p_el_debiased > 0.8
        OR z.p_cs_debiased > 0.8
        OR z.p_mg > 0.4
    )
\end{minted}

Il CSV locale non conserva \texttt{clean}, \texttt{type} e \(g\).
La verifica retrospettiva di questi vincoli richiede pertanto
l'accesso al catalogo sorgente, non il solo manifest finale.

\chapter{Codice del confronto appaiato}
\label{app:codice-appaiamento}

I frammenti seguenti provengono da \texttt{paired\_model\_comparison}
nel modulo \texttt{evaluation/statistics.py}. Sono estratti, non
programmi autonomi: import e controlli preliminari sono omessi e i
ritorni a capo sono adattati alla pagina. Il chiamante filtra prima
task e numero di qubit e passa Q come \texttt{model\_a}, B2 come
\texttt{model\_b}. Il confronto primario riguarda il solo Full;
i confronti secondari sono separati per numerosità.

\begin{code}
\begin{minted}{python}
pairing_keys = [
    key for key in (seed_column, "fold", "repeat", "subset_size")
    if key in frame.columns
]
selected = frame.loc[
    frame["model"].isin([model_a, model_b]),
    pairing_keys + ["model", metric]
]
if selected.duplicated(pairing_keys + ["model"]).any():
    raise ValueError("Esistono più risultati per la stessa coppia modello/fold/seed")
paired = selected.pivot(
    index=pairing_keys, columns="model", values=metric
).dropna()
\end{minted}
\caption{Costruzione delle coppie e rifiuto delle osservazioni duplicate.}
\label{lst:paired-code}
\end{code}

Il nome tecnico \texttt{fold} è conservato negli artefatti, ma nella
matrice analizzata vale sempre 1: non indica l'esecuzione di una
cross-validation completa. La funzione verifica inoltre che esistano
coppie complete prima di procedere. Nel \cref{lst:paired-bootstrap}
le differenze sono raggruppate per seed; ciascun contrasto della tesi
comprende dieci unità di ricampionamento, non 640 osservazioni.

\begin{code}
\begin{minted}{python}
paired_differences = (
    (paired[model_a] - paired[model_b])
    .rename("difference").reset_index()
)
inference_values = (
    paired_differences.groupby(seed_column, sort=True)["difference"]
    .mean().to_numpy(dtype=np.float64)
)
rng = np.random.default_rng(seed)
bootstrap_indices = rng.integers(
    0, len(inference_values),
    size=(n_bootstrap, len(inference_values))
)
bootstrap_means = inference_values[bootstrap_indices].mean(axis=1)
confidence_interval = np.quantile(bootstrap_means, [0.025, 0.975])
\end{minted}
\caption{Bootstrap percentile sulla media delle differenze appaiate.}
\label{lst:paired-bootstrap}
\end{code}

Il ricampionamento conserva l'appaiamento Q--B2, ma non corregge
la dipendenza tra holdout sovrapposti. Valgono quindi le limitazioni
discusse nei capitoli metodologico e di discussione. I p-value e la
correzione di Holm costituiscono passaggi distinti da questi intervalli.
